\documentclass[10pt,a4paper]{amsart}
\usepackage[margin=19mm]{geometry}
\usepackage{amsmath,amssymb,mathtools}
\usepackage[scr=rsfs,cal=euler]{mathalfa}
\usepackage{xfrac}
\usepackage[unicode,hidelinks,bookmarks=false]{hyperref}
\newcommand{\ie}{i.e.\ }
\newcommand{\cf}{cf.\ }
\newcommand{\resp}{respectively\ }
\newcommand{\Subsection}{\subsection}
\providecommand{\nobreakdash}{\nobreak\hbox{-}}
\setlength{\emergencystretch}{2em}
\multlinegap0pt
\usepackage{siunitx}
\usepackage{siunitx}
\usepackage{siunitx}
\usepackage{siunitx}
\usepackage[shortlabels]{enumitem}
\usepackage{bm}
\usepackage{amssymb}
\usepackage{siunitx}
\usepackage{xcolor}
\usepackage{tcolorbox}
\usepackage{mdframed}
\usepackage{tikz}
\usepackage{algorithm}\usepackage{algpseudocode}
\newtheorem{lemma}{Lemma}
\title{Dimension-Free Success Bounds for Constrained Quantum Optimization via Fejér Filtering}
\author{Chinonso Onah, Kristel Michielsen}
\date{Source: arXiv:2603.01809v3 (CC BY 4.0)}
\begin{document}
\maketitle

\noindent\emph{Source and licence.} Dimension-Free Success Bounds for Constrained Quantum Optimization via Fejér Filtering. Version arXiv:2603.01809v3 (CC BY 4.0), distributed by arXiv under the Creative Commons Attribution 4.0 International licence, http://creativecommons.org/licenses/by/4.0/, as recorded in the arXiv licence metadata for this exact version and shown on the abstract page https://arxiv.org/abs/2603.01809v3. Full licence notice: \texttt{licenses/arxiv-2603.01809-CC-BY-4.0.txt}.\par\smallskip

\noindent\textit{Editorial scope.} Counted unit: the article's lemma lem:lip-reduced (source0.tex lines 2868--2963). The article's complete original proof of it is delivered with it (source0.tex lines 2965--3098, one \texttt{proof} environment of 134 lines). No mathematics is written, repaired or re-derived: the statement and the proof are the article's own lines, in the article's own notation, and no retained line is substituted or re-flowed.  No further source-local context is needed: every cross-reference of every retained line resolves inside the two delivered spans.  Nothing was omitted from any retained span, so the package carries no omission record.  No retained line contains a citation, so the wrapper carries no bibliography and no bibliography entry.  No retained line contains a figure environment, an image-inclusion command or drawing code, so no image is delivered and the package declares no asset dependency.  Editorial formatting only, in addition: an AMS \texttt{amsart} preamble that restates the article's own declarations and macros used by the retained lines, namely the environments \texttt{lemma} on separate counters, together with the standard packages those lines need.  The retained environments print in this document as Lemma 1 in order of appearance, whereas the article prints the same items with the numbering of its own full document.  The complete native source of the article as distributed in the arXiv e-print for this exact version is bundled unchanged as \texttt{sources/195578/source0.tex}.
\begin{lemma}[Lipschitz and main--lobe control at reduced order]
\label{lem:lip-reduced}
Let \(p'=p-k\) with \(k\in\{1,\dots,p-1\}\). Then:
\begin{enumerate}[leftmargin=1.5em]

\item \textbf{Mixer-envelope mass on the optimum set.}
Define
\[
C_{\beta}^{(p')}
:=
\sum_{z\in\Omega^\star}
W_{p'}(z;\boldsymbol\beta).
\]
For any \(\boldsymbol\beta\) and perturbation
\(\Delta\boldsymbol\beta\),
\[
\left|
C_{\beta+\Delta\beta}^{(p')}
-
C_{\beta}^{(p')}
\right|
\le
L_W^{(p')}
\|\Delta\boldsymbol\beta\|_\infty,
\qquad
L_W^{(p')}
=
O\!\bigl(p'\|H_M\|\bigr).
\]

\item \textbf{Fej\'er main lobe (constant-fraction capture).}
The Fej\'er kernel
\[
F_{p'}(\Delta)
=
\frac{1}{p'+1}
\frac{
\sin^2\!\bigl((p'+1)\Delta/2\bigr)
}{
\sin^2(\Delta/2)
}
\]
has its nearest nonzero zeros at
\[
|\Delta|
=
\frac{2\pi}{p'+1}
\]
and hence has main-lobe width \(\Theta(1/p')\). For every fixed
\(c\in(0,\pi)\),
\[
|\Delta\theta|
\le
\frac{c}{p'+1}
\quad\Longrightarrow\quad
F_{p'}(\Delta\theta)
\ge
\kappa_c(p'+1),
\]
where
\[
\kappa_c
:=
\left(
\frac{\sin(c/2)}{c/2}
\right)^2
\in(0,1).
\]

Equivalently, let
\[
\Delta\theta
=
\Delta\gamma\,
\bigl(E(z)-E^\star\bigr)
\]
and choose
\[
R
\ge
\max_z
\bigl|E(z)-E^\star\bigr|.
\]
Then
\[
|\Delta\gamma|
\le
\frac{c}{(p'+1)R}
\quad\Longrightarrow\quad
F_{p'}(\Delta\theta)
\ge
\kappa_c(p'+1).
\]

\end{enumerate}
\end{lemma}

\begin{proof}[Proof sketch]
\emph{(1) Mixer-envelope mass.}
Write
\[
U_M^{(p')}(\boldsymbol\beta)
=
\prod_{\ell=1}^{p'}
e^{-i\beta_\ell H_M}
\]
and introduce the projector onto the full optimum subspace,
\[
\Pi_{\Omega^\star}
:=
\sum_{z\in\Omega^\star}
|z\rangle\langle z|.
\]
Then
\[
C_{\beta}^{(p')}
=
\langle s|
U_M^{(p')}(\boldsymbol\beta)^\dagger
\Pi_{\Omega^\star}
U_M^{(p')}(\boldsymbol\beta)
|s\rangle.
\]
By the Wilcox--Duhamel formula,
\[
\left\|
\frac{\partial}{\partial\beta_\ell}
e^{-i\beta_\ell H_M}
\right\|
\le
\|H_M\|.
\]
A product rule together with unitarity gives
\[
\left\|
\frac{\partial}{\partial\beta_\ell}
U_M^{(p')}(\boldsymbol\beta)
\right\|
\le
\|H_M\|.
\]
Since \(\Pi_{\Omega^\star}\) is an orthogonal projector,
\[
\|\Pi_{\Omega^\star}\|=1,
\]
independently of the degeneracy of the optimum. Therefore,
\[
\left|
C_{\beta+\Delta\beta}^{(p')}
-
C_{\beta}^{(p')}
\right|
\le
2
\sum_{\ell=1}^{p'}
\left\|
\frac{\partial}{\partial\beta_\ell}
U_M^{(p')}(\widetilde{\boldsymbol\beta})
\right\|
|\Delta\beta_\ell|
\le
2p'\|H_M\|
\|\Delta\boldsymbol\beta\|_\infty,
\]
for some \(\widetilde{\boldsymbol\beta}\) on the segment joining
\(\boldsymbol\beta\) and
\(\boldsymbol\beta+\Delta\boldsymbol\beta\). Hence
\[
L_W^{(p')}
=
O\!\bigl(p'\|H_M\|\bigr).
\]

\smallskip
\emph{(2) Fej\'er main lobe.}
For
\[
|\Delta|
\le
\frac{c}{p'+1},
\]
write \(u=(p'+1)\Delta/2\). Then \(|u|\le c/2<\pi/2\), and
\[
F_{p'}(\Delta)
=
(p'+1)
\left(
\frac{\sin u}{u}
\right)^2
\left(
\frac{\Delta/2}{\sin(\Delta/2)}
\right)^2.
\]
The standard \(\sin x/x\) bounds give
\[
F_{p'}(\Delta)
\ge
(p'+1)
\left(
\frac{\sin(c/2)}{c/2}
\right)^2
=
\kappa_c(p'+1).
\]
Finally, if
\[
\Delta\theta
=
\Delta\gamma\,
\bigl(E(z)-E^\star\bigr)
\]
and
\[
R
\ge
\max_z|E(z)-E^\star|,
\]
then
\[
|\Delta\gamma|
\le
\frac{c}{(p'+1)R}
\]
implies
\[
|\Delta\theta|
\le
\frac{c}{p'+1},
\]
which proves the stated phase-resolution bound.
\end{proof}

\end{document}
